\section{模型假设}

% 写作提示：只保留模型确实需要且后文会使用的假设；每条说明依据、简化作用、适用范围和可能影响。
% 下方假设仅为示例，不得无依据地照搬”数据完全可靠、因素独立、系统稳定、变量连续可微”等表述。

为简化问题，本文做出以下基本假设：
\begin{enumerate}
    \item \textbf{模态独立性假设}：假设文本、语音、视觉三种模态在特征层面相对独立，模态间的交互可以通过集成学习的隐式建模捕获，无需显式设计复杂的跨模态注意力机制。该假设简化了模型结构，适用于中等规模数据集，但可能忽略强相关场景（如语音韵律与面部表情的同步性）。

    \item \textbf{时序平稳性假设}：假设情感在单个样本内部的时序演化可以通过统计特征（均值、方差等）充分表征，无需建模精细的时序依赖关系。该假设支持统计聚合方法的使用，但对情感突变（如情绪爆发）的捕捉能力较弱。

    \item \textbf{缺失随机性假设}：假设模态缺失机制为随机缺失（MAR, Missing At Random），即缺失概率与观测值有关但与未观测值无关。该假设支持训练时随机Dropout策略的有效性，但若缺失与情感标签本身强相关（如极端情绪导致采集失败），模型泛化性能可能下降。

    \item \textbf{标签连续性假设}：假设情感标签在[-3, 3]区间内连续分布，可以采用回归模型而非离散分类。该假设符合心理学中情感维度理论（效价-唤醒模型），但对边界情感（如极度悲伤-3或极度喜悦+3）的预测精度可能受限。

    \item \textbf{特征充分性假设}：假设BERT文本编码、COVAREP音频特征和面部动作单元能够充分表征情感相关信息，无需引入额外的外部知识或上下文。该假设支持端到端建模，但对需要常识推理的情感识别（如讽刺、隐喻）可能效果不佳。
\end{enumerate}

